\epigraph{The notion of a complete market is the central concept in the theory of allocation under uncertainty.}{\textit{Kenneth Arrow}}


This chapter develops the theory of dynamic portfolio choice, building on the static CES optimization framework. We begin with Epstein-Zin preferences, which separate risk aversion from intertemporal substitution, then show how the standard time-separable case emerges as a special restriction. The key insight is that portfolio choice is a nested problem inside the consumption-savings decision, and homothetic preferences deliver a powerful separation theorem.

% Section 1: Static Portfolio Choice Revisited
% ==============================================================
% Section 1: Static Portfolio Choice Revisited
% ==============================================================



\section{Static Portfolio Choice Revisited}
\label{sec:static-portfolio}

This section connects the static CES optimization framework developed earlier to the portfolio choice problem. We begin with complete markets, where the closed-form solutions follow immediately from our canonical problem, then examine how incomplete markets break this tractability.

% --------------------------------------------------------------
\subsection{The Static Problem}
\label{subsec:static-problem}
% --------------------------------------------------------------

Consider an investor with wealth $W$ facing $S$ states of the world. State $s$ occurs with probability $\pi_s > 0$, where $\sum_{s=1}^{S} \pi_s = 1$. The investor chooses state-contingent consumption $(c_1, \ldots, c_S)$ to maximize a certainty equivalent.

\paragraph{CRRA Preferences.}

With constant relative risk aversion $\gamma > 0$, expected utility is
\begin{equation}
\E[u(c)] = \sum_{s=1}^{S} \pi_s \frac{c_s^{1-\gamma}}{1-\gamma}.
\label{eq:expected-utility}
\end{equation}
The certainty equivalent---the sure consumption yielding the same utility as the lottery---is
\begin{equation}
\CE(\bc; \bm{\pi}, \gamma) = \left( \sum_{s=1}^{S} \pi_s \, c_s^{1-\gamma} \right)^{\frac{1}{1-\gamma}}.
\label{eq:certainty-equivalent}
\end{equation}

\paragraph{Connection to Generalized Means.}

The certainty equivalent is a \emph{classical generalized mean} with power parameter $\rho = 1 - \gamma$ and weights equal to probabilities:
\begin{equation}
\CE(\bc; \bm{\pi}, \gamma) = \Mbar(\bc; \bm{\pi}, 1-\gamma).
\label{eq:ce-as-mean}
\end{equation}
Risk aversion $\gamma$ maps to the power parameter via $\rho = 1 - \gamma$:
\begin{center}
\begin{tabular}{@{}lccc@{}}
\toprule
Risk Aversion & $\gamma < 1$ & $\gamma = 1$ & $\gamma > 1$ \\
\midrule
Power $\rho = 1-\gamma$ & $\rho > 0$ & $\rho = 0$ & $\rho < 0$ \\
Mean Type & Above geometric & Geometric & Below geometric \\
Risk Attitude & Low aversion & Log utility & High aversion \\
\bottomrule
\end{tabular}
\end{center}

Higher risk aversion ($\gamma > 1$) corresponds to $\rho < 0$, which pulls the certainty equivalent below the arithmetic mean---the investor penalizes dispersion across states.

% --------------------------------------------------------------
\subsection{Complete Markets}
\label{subsec:complete-markets}
% --------------------------------------------------------------

\paragraph{Arrow-Debreu Securities.}

A market is \emph{complete} if there exist $S$ linearly independent assets spanning all states. The canonical case is $S$ Arrow-Debreu securities, where security $s$ pays one unit if state $s$ occurs and zero otherwise. Let $q_s > 0$ denote the price of the Arrow-Debreu security for state $s$.

The investor's problem is
\begin{equation}
\max_{c_1, \ldots, c_S} \; \CE(\bc; \bm{\pi}, \gamma) \quad \text{subject to} \quad \sum_{s=1}^{S} q_s c_s = W.
\label{eq:complete-markets-problem}
\end{equation}

\paragraph{This is the Canonical Problem.}

Problem \eqref{eq:complete-markets-problem} is exactly the CES optimization problem from Chapter~\ref{ch:canonical-problems}:
\begin{itemize}[nosep]
    \item Objective: generalized mean $\Mbar(\bc; \bm{\pi}, 1-\gamma)$
    \item Weights: probabilities $\pi_s$
    \item Prices: Arrow-Debreu prices $q_s$
    \item Budget: initial wealth $W$
\end{itemize}

Applying our earlier results immediately yields the solution.

\begin{proposition}[Complete Markets Solution]
\label{prop:complete-markets}
With complete Arrow-Debreu markets, optimal state-contingent consumption is
\begin{equation}
c_s\opt = \pi_s^{1/\gamma} \left( \frac{q_s}{\Qstar} \right)^{-1/\gamma} \cdot \frac{W}{\Qstar},
\label{eq:complete-markets-demand}
\end{equation}
where the \emph{risk-adjusted price index} is
\begin{equation}
\Qstar = \left( \sum_{s=1}^{S} \pi_s^{1/\gamma} \, q_s^{1-1/\gamma} \right)^{\frac{\gamma}{\gamma-1}}.
\label{eq:risk-price-index}
\end{equation}
The maximized certainty equivalent is
\begin{equation}
\CE\opt = \frac{W}{\Qstar}.
\label{eq:ce-value}
\end{equation}
\end{proposition}

\begin{proof}
Apply Proposition~\ref{prop:value} with $\omega_s = \pi_s$, $p_s = q_s$, $\rho = 1 - \gamma$, and $B = W$. The elasticity of substitution across states is $\varepsilon = 1/(1 - \rho) = 1/\gamma$.
\end{proof}

\paragraph{Interpretation.}

The demand function \eqref{eq:complete-markets-demand} has the standard multiplicative structure:
\[
c_s\opt = \underbrace{\pi_s^{1/\gamma}}_{\text{probability}} \times \underbrace{\left( \frac{q_s}{\Qstar} \right)^{-1/\gamma}}_{\text{relative price}} \times \underbrace{\frac{W}{\Qstar}}_{\text{real wealth}}.
\]
States with higher probability receive more consumption. States with higher Arrow-Debreu prices receive less consumption, with the response governed by the risk tolerance $1/\gamma$. The real wealth $W/\Qstar$ scales all consumption proportionally.

\begin{calloutnote}[State Prices and Risk-Neutral Probabilities]
Define the \emph{risk-neutral probability} $\tilde{\pi}_s \equiv q_s / \sum_j q_j$. Then
\[
\frac{q_s}{\pi_s} = \frac{\tilde{\pi}_s}{\pi_s} \cdot \sum_j q_j
\]
is the \emph{pricing kernel} or \emph{stochastic discount factor} for state $s$, up to a constant. States where the pricing kernel is high (consumption is expensive) receive less consumption. This is the foundation of consumption-based asset pricing.
\end{calloutnote}

% --------------------------------------------------------------
\subsection{Incomplete Markets}
\label{subsec:incomplete-markets}
% --------------------------------------------------------------

Now suppose only $J < S$ assets are available for trade. Let $R_s^j$ denote the gross return of asset $j$ in state $s$. If the investor allocates portfolio share $\alpha_j$ to asset $j$ (with $\sum_{j=1}^{J} \alpha_j = 1$), wealth in state $s$ is
\begin{equation}
W_s = W \sum_{j=1}^{J} \alpha_j R_s^j.
\label{eq:portfolio-wealth}
\end{equation}

The investor's problem becomes
\begin{equation}
\max_{\alpha_1, \ldots, \alpha_J} \; \CE\left( W \sum_{j} \alpha_j R_1^j, \ldots, W \sum_{j} \alpha_j R_S^j \right).
\label{eq:incomplete-markets-problem}
\end{equation}

\paragraph{The Spanning Constraint.}

The key difference from complete markets is that consumption across states is no longer freely chosen. Instead, the consumption vector $\bc = (c_1, \ldots, c_S)$ must lie in the \emph{asset span}---the $J$-dimensional subspace of $\mathbb{R}^S$ generated by the columns of the return matrix:
\begin{equation}
\bc \in \text{span}\{ \bR^1, \ldots, \bR^J \}, \quad \text{where } \bR^j = (R_1^j, \ldots, R_S^j)'.
\label{eq:asset-span}
\end{equation}

\paragraph{Why Closed Forms Fail.}

With complete markets, we optimize over all of $\mathbb{R}^S_+$---there are $S$ choice variables and $S$ first-order conditions, yielding a ``square'' system with closed-form solutions.

With incomplete markets, we have only $J < S$ choice variables (portfolio weights), but the first-order conditions involve expectations over $S$ states. The certainty equivalent
\begin{equation}
\CE = \left( \sum_{s=1}^{S} \pi_s \left( W \sum_{j=1}^{J} \alpha_j R_s^j \right)^{1-\gamma} \right)^{\frac{1}{1-\gamma}}
\label{eq:ce-incomplete}
\end{equation}
is a nonlinear function of the portfolio weights $\balpha$. Differentiating with respect to $\alpha_j$:
\begin{equation}
\frac{\partial \CE}{\partial \alpha_j} = \CE^{\gamma} \cdot W \sum_{s=1}^{S} \pi_s \, W_s^{-\gamma} \, R_s^j.
\label{eq:portfolio-foc}
\end{equation}

The first-order condition $\partial \CE / \partial \alpha_j = \lambda$ (for all $j$) requires
\begin{equation}
\E\left[ W_s^{-\gamma} R_s^j \right] = \text{constant for all } j.
\label{eq:euler-static}
\end{equation}

This is a system of $J$ equations in $J$ unknowns, but the expectation $\E[W_s^{-\gamma} R_s^j]$ depends on the \emph{entire distribution} of portfolio returns, not just moments. The nonlinear interaction between $W_s^{-\gamma}$ and $R_s^j$ generically precludes closed-form solutions.

\begin{calloutnote}[Why No Closed Form?]
The complete-markets solution $c_s\opt \propto \pi_s^{1/\gamma} q_s^{-1/\gamma}$ generically does not lie in the asset span. The incomplete-markets solution is the ``closest'' feasible point to this ideal, where distance is measured in the metric induced by the certainty equivalent. This projection onto a lower-dimensional subspace has no closed form except in special cases.

Geometrically: we seek the point in a $J$-dimensional affine subspace that maximizes a strictly concave objective over $\mathbb{R}^S$. The solution exists and is unique, but characterizing it requires solving a nonlinear system.
\end{calloutnote}

\begin{calloutnote}[Special Cases with Closed Forms]
Analytical solutions exist in three important cases:
\begin{enumerate}[label=(\roman*), nosep]
    \item \textbf{Log utility} ($\gamma = 1$): The objective becomes $\E[\log W_s]$, which separates additively. Portfolio choice is myopic and depends only on expected log returns.
    
    \item \textbf{CARA utility with normal returns}: With exponential utility $u(W) = -e^{-aW}$ and jointly normal returns, the problem reduces to mean-variance optimization, yielding the classic two-fund separation.
    
    \item \textbf{CRRA utility with log-normal returns}: When log returns are normal, approximate closed forms exist via log-linearization (Campbell-Viceira). The next subsection develops the exact log-utility case.
\end{enumerate}
Outside these cases, numerical methods are required---a theme developed in later chapters.
\end{calloutnote}

\paragraph{Other Linear Constraints.}

The incomplete markets problem is one instance of a broader class: CES optimization subject to linear constraints on the choice vector. Several other economically important constraints share this structure:
\begin{itemize}[nosep]
    \item \textbf{Short-sale constraints}: The requirement $\alpha_j \geq 0$ restricts portfolios to a cone within the asset span. Binding constraints create kinks in the value function.
    
    \item \textbf{Leverage constraints}: A bound $\sum_j |\alpha_j| \leq L$ or $\alpha_j \geq -\bar{\alpha}$ limits borrowing, truncating the feasible region.
    
    \item \textbf{Regulatory constraints}: Capital requirements, position limits, or concentration limits impose linear inequalities on portfolio weights.
    
    \item \textbf{Linear transaction costs}: Proportional costs $\kappa |\text{trade}_j|$ can be reformulated as linear constraints on rebalancing, though they introduce state-dependence.
\end{itemize}

In each case, the mathematical structure is identical: maximize a CES objective over a polyhedron---the intersection of half-spaces---rather than the full asset span. The complete-markets closed form provides the \emph{unconstrained} optimum; the actual solution is a projection onto the feasible set. When constraints bind, Lagrange multipliers have natural interpretations as shadow prices of the restricted resources.

% --------------------------------------------------------------
\subsection{The Log-Normal / Log-Utility Case}
\label{subsec:log-utility}
% --------------------------------------------------------------

The log-utility case ($\gamma = 1$) with log-normal returns provides a complete analytical solution and serves as the tractability benchmark.

\paragraph{Setup.}

Consider two assets: a risk-free bond with gross return $R^f$ and a risky asset with random gross return $R$. Assume log returns are normal:
\begin{equation}
\log R \sim N(\mu - \tfrac{1}{2}\sigma^2, \sigma^2),
\label{eq:log-return}
\end{equation}
where the parametrization ensures $\E[R] = e^{\mu}$. Let $\alpha$ denote the portfolio share in the risky asset.

\paragraph{The Log-Utility Objective.}

With log utility, the certainty equivalent is the geometric mean:
\begin{equation}
\CE = \exp\left( \E[\log W_1] \right) = \exp\left( \log W_0 + \E[\log R^p] \right),
\label{eq:log-ce}
\end{equation}
where $R^p = (1-\alpha) R^f + \alpha R$ is the portfolio return. Maximizing $\CE$ is equivalent to maximizing expected log return:
\begin{equation}
\max_{\alpha} \; \E\left[ \log\left( (1-\alpha) R^f + \alpha R \right) \right].
\label{eq:log-utility-problem}
\end{equation}

\paragraph{The Solution.}

For log-normal returns, the first-order condition yields the celebrated result.

\begin{proposition}[Log-Utility Portfolio]
\label{prop:log-portfolio}
With log utility and a risky asset satisfying \eqref{eq:log-return}, the optimal portfolio share is
\begin{equation}
\alpha\opt = \frac{\mu - \log R^f}{\sigma^2} = \frac{\E[\log R] + \frac{1}{2}\sigma^2 - \log R^f}{\sigma^2}.
\label{eq:log-portfolio}
\end{equation}
Equivalently, using the arithmetic risk premium $\E[R] - R^f \approx \mu - \log R^f + \frac{1}{2}\sigma^2$:
\begin{equation}
\alpha\opt \approx \frac{\E[R] - R^f}{\sigma^2}.
\label{eq:log-portfolio-approx}
\end{equation}
\end{proposition}

This is the \emph{Kelly criterion}---the growth-optimal portfolio that maximizes the long-run geometric growth rate of wealth.

\paragraph{Key Properties.}

The log-utility solution has remarkable features:

\begin{enumerate}[label=(\roman*)]
    \item \textbf{Myopic}: The optimal portfolio share depends only on the current-period return distribution, not on the investment horizon or future opportunities.
    
    \item \textbf{Wealth-independent}: The share $\alpha\opt$ does not depend on wealth $W_0$. Rich and poor investors hold the same portfolio proportions.
    
    \item \textbf{Separation}: Portfolio choice is independent of the consumption-savings decision. We can solve for $\alpha\opt$ first, then determine optimal consumption.
\end{enumerate}

\begin{calloutimportant}[Why Log Utility is Special]
Log utility ($\gamma = 1$) corresponds to $\rho = 0$ in the generalized mean---the geometric mean. This is the boundary between means that favor high values ($\rho > 0$) and means that favor low values ($\rho < 0$).

At this boundary, the certainty equivalent has a multiplicative structure:
\[
\CE = \left( \prod_{s} c_s^{\pi_s} \right) = \exp\left( \sum_s \pi_s \log c_s \right).
\]
The logarithm converts the product to a sum, making the objective \emph{additively separable} across states. This separability is what delivers closed forms.

For $\gamma \neq 1$, the power function $c^{1-\gamma}$ does not separate, and the nonlinear interaction between states prevents analytical solutions.
\end{calloutimportant}

\paragraph{Beyond Log Utility.}

For general $\gamma \neq 1$, the optimal portfolio share with log-normal returns satisfies
\begin{equation}
\alpha\opt \approx \frac{1}{\gamma} \cdot \frac{\mu - \log R^f}{\sigma^2}.
\label{eq:crra-portfolio-approx}
\end{equation}
This approximation (Merton, Campbell-Viceira) shows that higher risk aversion $\gamma$ reduces the risky share proportionally. However, the exact solution requires numerical methods because $\E[W^{1-\gamma}]$ depends on the full distribution, not just $\mu$ and $\sigma^2$.

The factor $1/\gamma$ is the \emph{risk tolerance}. It measures the investor's willingness to bear risk in exchange for higher expected returns. Log utility ($\gamma = 1$) has unit risk tolerance; higher risk aversion reduces the risky allocation.

\begin{callouttip}[The Merton Share]
Equation \eqref{eq:crra-portfolio-approx} is often called the \emph{Merton share} after Merton (1969). In continuous time with geometric Brownian motion, this expression is exact. The discrete-time/log-normal case is an approximation that becomes exact as the time interval shrinks.
\end{callouttip}



% Section 2: Epstein-Zin Preferences
% ==============================================================
% Section 2: Epstein-Zin Preferences
% ==============================================================

\section{Epstein-Zin Preferences}
\label{sec:epstein-zin}

The standard time-separable expected utility model ties together two conceptually distinct aspects of preferences: aversion to risk (across states) and aversion to intertemporal fluctuations (across time). Epstein-Zin preferences untie this knot, allowing risk aversion and the elasticity of intertemporal substitution to be specified independently. This separation is not merely a generalization for its own sake---it resolves empirical puzzles and provides a cleaner conceptual framework for dynamic portfolio choice.

% --------------------------------------------------------------
\subsection{The Recursive Aggregator}
\label{subsec:ez-aggregator}
% --------------------------------------------------------------

\paragraph{Motivation.}

Consider two distinct questions an investor might face:
\begin{enumerate}[nosep]
    \item \textbf{Risk aversion}: How much would you pay to avoid a fair gamble over consumption \emph{today}?
    \item \textbf{Intertemporal substitution}: How willing are you to shift consumption between \emph{today and tomorrow} in response to interest rate changes?
\end{enumerate}

These are logically independent questions. An investor might be highly risk averse (unwilling to bear uncertainty) yet also highly willing to substitute intertemporally (responsive to interest rates). Or vice versa.

With standard time-separable CRRA utility,
\begin{equation}
U = \E_0 \left[ \sum_{t=0}^{\infty} \beta^t \frac{c_t^{1-\gamma}}{1-\gamma} \right],
\label{eq:time-separable}
\end{equation}
both questions are answered by the \emph{same} parameter $\gamma$: risk aversion equals $\gamma$, and the elasticity of intertemporal substitution (IES) equals $1/\gamma$. This is a strong restriction with no economic justification.

\paragraph{The Epstein-Zin Aggregator.}

Epstein and Zin (1989) proposed a recursive utility specification that separates these two roles. The value function $V_t$ satisfies
\begin{equation}
\boxed{V_t = \left[ (1-\beta) \, c_t^{1-1/\theta} + \beta \left( \E_t\left[ V_{t+1}^{1-\gamma} \right] \right)^{\frac{1-1/\theta}{1-\gamma}} \right]^{\frac{1}{1-1/\theta}},}
\label{eq:epstein-zin}
\end{equation}
where:
\begin{itemize}[nosep]
    \item $\theta > 0$ is the \textbf{elasticity of intertemporal substitution} (IES)
    \item $\gamma > 0$ is the \textbf{coefficient of relative risk aversion}
    \item $\beta \in (0,1)$ is the \textbf{discount factor}
\end{itemize}

The recursion has two nested aggregations, each with CES structure.

\paragraph{Two Nested CES Aggregations.}

Equation \eqref{eq:epstein-zin} combines two generalized means:

\begin{enumerate}[label=(\roman*)]
    \item \textbf{Aggregation across states} (inner): The term
    \begin{equation}
    \mu_t(V_{t+1}) \equiv \left( \E_t\left[ V_{t+1}^{1-\gamma} \right] \right)^{\frac{1}{1-\gamma}}
    \label{eq:certainty-equivalent-v}
    \end{equation}
    is the \emph{certainty equivalent} of next period's value. This is a generalized mean with power $\rho_{\text{risk}} = 1 - \gamma$ and weights equal to probabilities. Risk aversion $\gamma$ governs this aggregation.
    
    \item \textbf{Aggregation across time} (outer): Given the certainty equivalent $\mu_t$, the recursion
    \begin{equation}
    V_t = \left[ (1-\beta) \, c_t^{1-1/\theta} + \beta \, \mu_t^{1-1/\theta} \right]^{\frac{1}{1-1/\theta}}
    \label{eq:intertemporal-aggregation}
    \end{equation}
    is a CES aggregator over $(c_t, \mu_t)$ with power $\rho_{\text{time}} = 1 - 1/\theta$. The IES $\theta$ governs this aggregation.
\end{enumerate}

\begin{calloutimportant}[The Two-Layer Structure]
Epstein-Zin preferences nest two CES aggregations:
\begin{center}
\begin{tabular}{@{}lcc@{}}
\toprule
\textbf{Aggregation} & \textbf{Power Parameter} & \textbf{Elasticity} \\
\midrule
Across states (risk) & $\rho_{\text{risk}} = 1 - \gamma$ & $1/\gamma$ \\
Across time (intertemporal) & $\rho_{\text{time}} = 1 - 1/\theta$ & $\theta$ \\
\bottomrule
\end{tabular}
\end{center}
The certainty equivalent $\mu_t$ serves as the ``price'' of future utility in the intertemporal aggregation. This nested structure is the key to tractability.
\end{calloutimportant}

\paragraph{Rewriting with Explicit Means.}

Using our notation for generalized means, \eqref{eq:epstein-zin} becomes
\begin{equation}
V_t = \M\left( c_t, \mu_t(V_{t+1}); (1-\beta, \beta), 1-\tfrac{1}{\theta} \right),
\label{eq:ez-as-mean}
\end{equation}
where
\begin{equation}
\mu_t(V_{t+1}) = \Mbar\left( V_{t+1}; \mathbf{p}_t, 1-\gamma \right)
\label{eq:mu-as-mean}
\end{equation}
and $\mathbf{p}_t$ is the vector of state probabilities conditional on information at $t$.

% --------------------------------------------------------------
\subsection{The Time-Separable Case}
\label{subsec:time-separable}
% --------------------------------------------------------------

\paragraph{The Restriction $\gamma = 1/\theta$.}

A remarkable simplification occurs when risk aversion equals the inverse of the IES:
\begin{equation}
\gamma = \frac{1}{\theta} \quad \Longleftrightarrow \quad \rho_{\text{risk}} = \rho_{\text{time}} = 1 - \gamma.
\label{eq:ez-restriction}
\end{equation}
Under this restriction, both aggregations use the \emph{same} power parameter. The nested structure collapses.

\begin{proposition}[Recovery of Time-Separable Utility]
\label{prop:time-separable}
When $\gamma = 1/\theta$, Epstein-Zin utility reduces to standard time-separable CRRA expected utility:
\begin{equation}
V_t = \left( \E_t \left[ \sum_{s=t}^{\infty} \beta^{s-t} c_s^{1-\gamma} \right] \right)^{\frac{1}{1-\gamma}}.
\label{eq:ez-to-crra}
\end{equation}
Equivalently, maximizing $V_t$ is equivalent to maximizing $\E_t[\sum_{s=t}^{\infty} \beta^{s-t} u(c_s)]$ with $u(c) = c^{1-\gamma}/(1-\gamma)$.
\end{proposition}

\begin{proof}
With $\gamma = 1/\theta$, write $\rho = 1 - \gamma = 1 - 1/\theta$. The recursion \eqref{eq:epstein-zin} becomes
\begin{align}
V_t &= \left[ (1-\beta) c_t^{\rho} + \beta \left( \E_t[V_{t+1}^{\rho}] \right)^{\rho/\rho} \right]^{1/\rho} \notag \\
&= \left[ (1-\beta) c_t^{\rho} + \beta \, \E_t[V_{t+1}^{\rho}] \right]^{1/\rho}.
\label{eq:collapsed-recursion}
\end{align}
Define $\tilde{V}_t = V_t^{\rho}$. Then
\[
\tilde{V}_t = (1-\beta) c_t^{\rho} + \beta \, \E_t[\tilde{V}_{t+1}].
\]
This is a linear recursion with solution
\[
\tilde{V}_t = (1-\beta) \E_t \left[ \sum_{s=t}^{\infty} \beta^{s-t} c_s^{\rho} \right].
\]
Taking the $1/\rho$ power and noting that monotonic transformations preserve optima yields \eqref{eq:ez-to-crra}.
\end{proof}

\paragraph{Why This Restriction is Problematic.}

The restriction $\gamma = 1/\theta$ implies:
\begin{itemize}[nosep]
    \item High risk aversion ($\gamma > 1$) $\Rightarrow$ low IES ($\theta < 1$): Risk-averse investors are also unwilling to substitute intertemporally.
    \item Low risk aversion ($\gamma < 1$) $\Rightarrow$ high IES ($\theta > 1$): Risk-tolerant investors are also highly responsive to interest rates.
\end{itemize}

Empirically, these implications are problematic:
\begin{itemize}[nosep]
    \item The equity premium puzzle suggests $\gamma \approx 10$ or higher to rationalize observed risk premia.
    \item But $\theta = 1/\gamma = 0.1$ implies extremely low intertemporal substitution, inconsistent with consumption responses to interest rate changes.
\end{itemize}

Epstein-Zin preferences resolve this by allowing, say, $\gamma = 10$ and $\theta = 1.5$ simultaneously---high risk aversion with moderate intertemporal substitution.

\begin{calloutnote}[The Unit Elasticity Benchmark]
When $\gamma = \theta = 1$, both aggregations become geometric means (Cobb-Douglas), and preferences reduce to
\[
V_t = c_t^{1-\beta} \cdot \exp\left( \beta \, \E_t[\log V_{t+1}] \right).
\]
Maximizing $\log V_t$ gives the additively separable objective $\E_t[\sum_s \beta^{s-t} \log c_s]$. This is the log-utility case from Section~\ref{subsec:log-utility}, which admits closed-form portfolio solutions.
\end{calloutnote}

% --------------------------------------------------------------
\subsection{Preference for Early or Late Resolution of Uncertainty}
\label{subsec:resolution-timing}
% --------------------------------------------------------------

A distinctive feature of Epstein-Zin preferences is that they encode attitudes toward the \emph{timing} of uncertainty resolution, not just its presence.

\paragraph{The Key Parameter.}

Define
\begin{equation}
\psi \equiv \frac{1 - \gamma}{1 - 1/\theta} = \frac{\rho_{\text{risk}}}{\rho_{\text{time}}}.
\label{eq:psi-definition}
\end{equation}
This ratio determines preferences over the timing of resolution:
\begin{itemize}[nosep]
    \item $\psi < 1$ (equivalently, $\gamma > 1/\theta$): \textbf{Preference for early resolution}. The agent prefers to learn outcomes sooner rather than later, even if no action can be taken.
    \item $\psi > 1$ (equivalently, $\gamma < 1/\theta$): \textbf{Preference for late resolution}. The agent prefers to delay learning outcomes.
    \item $\psi = 1$ (equivalently, $\gamma = 1/\theta$): \textbf{Indifference}. This is the time-separable case---timing of resolution is irrelevant.
\end{itemize}

\paragraph{Intuition.}

Consider two lotteries that deliver the same consumption streams with the same probabilities, but differ in when uncertainty is resolved:
\begin{itemize}[nosep]
    \item \textbf{Early resolution}: A coin is flipped today; you learn tomorrow's consumption now.
    \item \textbf{Late resolution}: The same coin flip occurs, but you learn the outcome tomorrow.
\end{itemize}

With time-separable utility, these are equivalent---only the probability distribution over final outcomes matters. With Epstein-Zin preferences and $\gamma > 1/\theta$, early resolution is preferred. The agent values ``knowing where they stand'' even before acting.

\begin{calloutimportant}[Why Timing Matters for Portfolio Choice]
Preference for early resolution ($\gamma > 1/\theta$) has important implications:
\begin{itemize}[nosep]
    \item Agents value assets that resolve uncertainty quickly (e.g., options, insurance).
    \item Long-horizon investors may behave differently than short-horizon investors \emph{even with identical terminal objectives}.
    \item The equity premium reflects not just risk aversion but also compensation for bearing persistent uncertainty.
\end{itemize}
This is one reason Epstein-Zin preferences have become central to modern asset pricing.
\end{calloutimportant}

% --------------------------------------------------------------
\subsection{The Ordinally Equivalent Formulation}
\label{subsec:ordinal-ez}
% --------------------------------------------------------------

For analytical work, it is often convenient to use an ordinally equivalent transformation of \eqref{eq:epstein-zin}.

\paragraph{The Transformed Value Function.}

Define $\tilde{V}_t = V_t^{1-\gamma}$ when $\gamma \neq 1$. The recursion becomes
\begin{equation}
\tilde{V}_t = \left[ (1-\beta) c_t^{1-1/\theta} + \beta \left( \E_t[\tilde{V}_{t+1}] \right)^{\frac{1-1/\theta}{1-\gamma}} \right]^{\frac{1-\gamma}{1-1/\theta}}.
\label{eq:ez-transformed}
\end{equation}

\paragraph{The Weil Parametrization.}

An alternative form, due to Weil (1990), writes
\begin{equation}
U_t = \left[ (1-\beta) c_t^{\rho} + \beta \left( \E_t[U_{t+1}^{\alpha}] \right)^{\rho/\alpha} \right]^{1/\rho},
\label{eq:weil-form}
\end{equation}
where $\rho = 1 - 1/\theta$ and $\alpha = (1-\gamma)/(1-1/\theta)$. This makes the nested CES structure explicit: power $\rho$ for intertemporal aggregation, power $\alpha \cdot \rho$ effectively for risk aggregation.

\paragraph{Special Cases.}

\begin{center}
\renewcommand{\arraystretch}{1.4}
\begin{tabular}{@{}lccl@{}}
\toprule
\textbf{Case} & $\gamma$ & $\theta$ & \textbf{Utility Form} \\
\midrule
Log utility & $1$ & $1$ & $\E[\sum_t \beta^t \log c_t]$ \\
Time-separable CRRA & $\gamma$ & $1/\gamma$ & $\E[\sum_t \beta^t c_t^{1-\gamma}/(1-\gamma)]$ \\
Risk-neutral, unit IES & $0$ & $1$ & $\E[\sum_t \beta^t c_t] / (1-\beta)$ \\
High RA, high IES & $10$ & $1.5$ & Full Epstein-Zin recursion \\
\bottomrule
\end{tabular}
\end{center}



% Section 3: The Consumption-Portfolio Problem
% ==============================================================
% Section 3: The Consumption-Portfolio Problem
% ==============================================================

\section{The Consumption-Portfolio Problem}
\label{sec:consumption-portfolio}

We now embed portfolio choice within the full dynamic consumption-savings problem. The key insight is that with homothetic preferences, the problem has a nested structure: portfolio allocation is an ``inner'' problem whose solution feeds into the ``outer'' consumption-savings decision. This separation dramatically simplifies the analysis.

% --------------------------------------------------------------
\subsection{The Sequential Problem}
\label{subsec:sequential-problem}
% --------------------------------------------------------------

\paragraph{Setup.}

Consider an investor with initial wealth $W_t$ at time $t$. The investor:
\begin{enumerate}[nosep]
    \item Chooses consumption $c_t$
    \item Allocates savings $W_t - c_t$ across $J$ assets with random returns $(R_{t+1}^1, \ldots, R_{t+1}^J)$
    \item Receives wealth $W_{t+1} = (W_t - c_t) \sum_{j=1}^{J} \alpha_j R_{t+1}^j$ next period
\end{enumerate}

Let $R_{t+1}^p(\balpha) = \sum_{j=1}^{J} \alpha_j R_{t+1}^j$ denote the portfolio return given weights $\balpha = (\alpha_1, \ldots, \alpha_J)$ with $\sum_j \alpha_j = 1$.

\paragraph{The Bellman Equation.}

With Epstein-Zin preferences, the value function satisfies
\begin{equation}
V_t(W_t) = \max_{c_t, \balpha} \left[ (1-\beta) c_t^{1-1/\theta} + \beta \left( \E_t\left[ V_{t+1}(W_{t+1})^{1-\gamma} \right] \right)^{\frac{1-1/\theta}{1-\gamma}} \right]^{\frac{1}{1-1/\theta}}
\label{eq:bellman-ez}
\end{equation}
subject to
\begin{equation}
W_{t+1} = (W_t - c_t) R_{t+1}^p(\balpha).
\label{eq:wealth-evolution}
\end{equation}

This is a joint optimization over consumption $c_t$ and portfolio weights $\balpha$. The challenge is that both decisions interact through the budget constraint and the continuation value.

% --------------------------------------------------------------
\subsection{Nesting of Decisions}
\label{subsec:nesting}
% --------------------------------------------------------------

The structure of \eqref{eq:bellman-ez} suggests a natural decomposition: first optimize the portfolio for any given level of savings, then optimize consumption taking into account the value of optimally invested savings.

\paragraph{The Inner Problem: Portfolio Choice.}

Fix savings $S_t = W_t - c_t > 0$. The portfolio problem is
\begin{equation}
\max_{\balpha} \; \mu_t\left( S_t \cdot R_{t+1}^p(\balpha) \right),
\label{eq:inner-problem}
\end{equation}
where $\mu_t(\cdot)$ denotes the certainty equivalent of next period's wealth mapped through the continuation value:
\begin{equation}
\mu_t(W_{t+1}) = \left( \E_t\left[ V_{t+1}(W_{t+1})^{1-\gamma} \right] \right)^{\frac{1}{1-\gamma}}.
\label{eq:mu-continuation}
\end{equation}

\paragraph{Homogeneity is the Key.}

The crucial observation is that with homothetic preferences, the value function is homogeneous in wealth:
\begin{equation}
V_{t+1}(W_{t+1}) = W_{t+1} \cdot v_{t+1},
\label{eq:value-homogeneity}
\end{equation}
where $v_{t+1}$ is the ``normalized'' value that depends on state variables other than wealth (e.g., investment opportunities, labor income ratios) but not on the level of wealth itself.

Substituting into \eqref{eq:mu-continuation}:
\begin{align}
\mu_t(S_t \cdot R_{t+1}^p) &= \left( \E_t\left[ \left( S_t \cdot R_{t+1}^p \cdot v_{t+1} \right)^{1-\gamma} \right] \right)^{\frac{1}{1-\gamma}} \notag \\
&= S_t \cdot \left( \E_t\left[ \left( R_{t+1}^p \cdot v_{t+1} \right)^{1-\gamma} \right] \right)^{\frac{1}{1-\gamma}}.
\label{eq:mu-factored}
\end{align}

The savings level $S_t$ factors out! This means the portfolio problem \eqref{eq:inner-problem} reduces to
\begin{equation}
\max_{\balpha} \; \left( \E_t\left[ \left( R_{t+1}^p(\balpha) \cdot v_{t+1} \right)^{1-\gamma} \right] \right)^{\frac{1}{1-\gamma}},
\label{eq:portfolio-problem-reduced}
\end{equation}
which is independent of the level of savings.

\begin{proposition}[Separation of Portfolio Choice]
\label{prop:portfolio-separation}
With homothetic Epstein-Zin preferences, the optimal portfolio weights $\balpha\opt$ are independent of:
\begin{enumerate}[label=(\roman*), nosep]
    \item Current wealth $W_t$
    \item Current consumption $c_t$
    \item The savings rate $s_t = (W_t - c_t)/W_t$
\end{enumerate}
Portfolio choice depends only on the distribution of returns and the normalized continuation value $v_{t+1}$.
\end{proposition}

\paragraph{The Outer Problem: Consumption-Savings.}

Given the optimal portfolio, define the \emph{certainty equivalent return}:
\begin{equation}
R_t^{CE} \equiv \left( \E_t\left[ \left( R_{t+1}^p(\balpha\opt) \cdot v_{t+1} \right)^{1-\gamma} \right] \right)^{\frac{1}{1-\gamma}}.
\label{eq:ce-return}
\end{equation}

The Bellman equation \eqref{eq:bellman-ez} becomes
\begin{equation}
V_t(W_t) = \max_{c_t} \left[ (1-\beta) c_t^{1-1/\theta} + \beta \left( (W_t - c_t) \cdot R_t^{CE} \right)^{1-1/\theta} \right]^{\frac{1}{1-1/\theta}}.
\label{eq:outer-problem}
\end{equation}

This is now a \emph{deterministic} optimization problem in $c_t$---all uncertainty has been absorbed into the certainty equivalent return $R_t^{CE}$.

\paragraph{The Consumption Rule.}

The first-order condition for \eqref{eq:outer-problem} yields the optimal consumption-wealth ratio.

\begin{proposition}[Optimal Consumption]
\label{prop:optimal-consumption}
The optimal consumption policy is linear in wealth:
\begin{equation}
c_t\opt = m_t \cdot W_t,
\label{eq:consumption-rule}
\end{equation}
where the marginal propensity to consume $m_t$ satisfies
\begin{equation}
\frac{m_t}{1 - m_t} = \left( \frac{1-\beta}{\beta} \right)^{\theta} \left( R_t^{CE} \right)^{\theta - 1}.
\label{eq:mpc-formula}
\end{equation}
Equivalently,
\begin{equation}
m_t = \frac{1}{1 + \left( \frac{\beta}{1-\beta} \right)^{\theta} \left( R_t^{CE} \right)^{\theta - 1}}.
\label{eq:mpc-explicit}
\end{equation}
\end{proposition}

\begin{proof}
Write $s = W_t - c_t$ for savings. The objective \eqref{eq:outer-problem} is
\[
\left[ (1-\beta) c^{1-1/\theta} + \beta (s \cdot R_t^{CE})^{1-1/\theta} \right]^{\frac{1}{1-1/\theta}}.
\]
The FOC with respect to $c$ (using $ds/dc = -1$) is
\[
(1-\beta) c^{-1/\theta} = \beta (s \cdot R_t^{CE})^{-1/\theta} \cdot R_t^{CE}.
\]
Rearranging:
\[
\frac{c}{s} = \left( \frac{1-\beta}{\beta} \right)^{\theta} (R_t^{CE})^{\theta - 1}.
\]
Since $c/s = m/(1-m)$ when $c = mW$ and $s = (1-m)W$, we obtain \eqref{eq:mpc-formula}.
\end{proof}

\begin{calloutnote}[The Role of the IES]
The elasticity $\theta$ governs how consumption responds to the certainty equivalent return:
\begin{itemize}[nosep]
    \item $\theta > 1$: Higher $R_t^{CE}$ \emph{reduces} current consumption (substitution effect dominates). The investor saves more to exploit good investment opportunities.
    \item $\theta < 1$: Higher $R_t^{CE}$ \emph{increases} current consumption (income effect dominates). The investor consumes more because future wealth will be higher.
    \item $\theta = 1$: Consumption-wealth ratio is constant regardless of investment opportunities (log utility).
\end{itemize}
This is the \emph{intertemporal} substitution effect, governed by $\theta$---completely separate from risk aversion $\gamma$.
\end{calloutnote}

% --------------------------------------------------------------
\subsection{The Certainty Equivalent of Wealth Growth}
\label{subsec:ce-wealth-growth}
% --------------------------------------------------------------

The certainty equivalent return $R_t^{CE}$ is the central object connecting portfolio choice to consumption-savings. It summarizes all relevant information about investment opportunities.

\paragraph{Interpretation.}

The certainty equivalent return answers the question: ``What risk-free return would make the investor indifferent between (i) investing optimally in risky assets and (ii) investing in a safe asset?''

If markets offered a risk-free return equal to $R_t^{CE}$, the investor would achieve the same utility as with the optimal risky portfolio.

\paragraph{Decomposition.}

When the normalized value $v_{t+1}$ is independent of returns (e.g., i.i.d. returns), we can decompose:
\begin{equation}
R_t^{CE} = \underbrace{\left( \E_t\left[ (R_{t+1}^{p,*})^{1-\gamma} \right] \right)^{\frac{1}{1-\gamma}}}_{\text{CE of portfolio return}} \times \underbrace{\left( \E_t\left[ v_{t+1}^{1-\gamma} \right] \right)^{\frac{1}{1-\gamma}}}_{\text{CE of continuation value}}.
\label{eq:ce-decomposition}
\end{equation}

The first term is the certainty equivalent of the optimally invested portfolio return---this is determined by the portfolio problem. The second term reflects uncertainty about future investment opportunities or other state variables.

\paragraph{The i.i.d. Case.}

When returns are i.i.d. and there are no other state variables, $v_{t+1} = v$ is constant. Then
\begin{equation}
R_t^{CE} = v \cdot \CE(R_{t+1}^{p,*}),
\label{eq:ce-iid}
\end{equation}
where $\CE(R^{p,*}) = (\E[(R^{p,*})^{1-\gamma}])^{1/(1-\gamma)}$ is the static certainty equivalent of the optimal portfolio return.

\begin{calloutimportant}[The Nested Structure]
The complete solution has a clean recursive structure:

\textbf{Step 1 (Portfolio):} Solve for optimal weights $\balpha\opt$ by maximizing the certainty equivalent of wealth growth. This is independent of current wealth and consumption.

\textbf{Step 2 (Consumption):} Given $\balpha\opt$, compute $R_t^{CE}$. Then solve for optimal consumption as a fraction of wealth using \eqref{eq:mpc-explicit}.

\textbf{Step 3 (Value):} The value function is $V_t(W_t) = v_t \cdot W_t$, where $v_t$ satisfies a recursion involving $R_t^{CE}$ and $m_t$.

This separation is the payoff from homothetic preferences: a high-dimensional joint optimization decomposes into tractable pieces.
\end{calloutimportant}

% --------------------------------------------------------------
\subsection{When Does Separation Fail?}
\label{subsec:separation-failure}
% --------------------------------------------------------------

The separation result relies on homotheticity. Several economically important features break this property.

\paragraph{Labor Income.}

If the investor receives non-tradeable labor income $Y_t$, the budget constraint becomes
\begin{equation}
W_{t+1} = (W_t + Y_t - c_t) R_{t+1}^p.
\label{eq:wealth-with-labor}
\end{equation}
Now savings $W_t + Y_t - c_t$ depends on the \emph{level} of wealth relative to income. The portfolio problem no longer separates because human capital (the present value of future labor income) is an implicit asset that cannot be traded.

\paragraph{Subsistence or Satiation.}

Preferences with a subsistence level $\bar{c}$,
\begin{equation}
u(c) = \frac{(c - \bar{c})^{1-\gamma}}{1-\gamma},
\label{eq:stone-geary}
\end{equation}
are not homothetic in $c$. The portfolio choice depends on how far consumption is from subsistence.

\paragraph{Borrowing Constraints.}

A constraint $W_{t+1} \geq \underline{W}$ introduces a kink in the value function. Near the constraint, the investor becomes effectively more risk averse, altering portfolio choice.

\paragraph{Background Risk.}

Uninsurable risks (health, unemployment, housing) interact with portfolio risk. Even if these risks are independent of asset returns, they affect the marginal utility of wealth and hence portfolio choice.

\begin{calloutnote}[Numerical Methods]
When separation fails, we must solve the joint consumption-portfolio problem numerically. The state space includes wealth plus any variables affecting labor income, constraints, or background risk. This is where the ``curse of dimensionality'' from Chapter~\ref{chap:introduction} bites---and where the computational methods developed later become essential.
\end{calloutnote}



% Section 4: The Value Function and Normalization
% ==============================================================
% Section 4: The Value Function and Normalization
% ==============================================================

\section{The Value Function and Normalization}
\label{sec:value-function}

The homothetic structure of Epstein-Zin preferences implies that the value function is proportional to wealth. This section exploits this property to derive the normalized Bellman equation, characterize the portfolio and consumption first-order conditions, and establish the Euler equations that govern intertemporal optimality.

% --------------------------------------------------------------
\subsection{Homogeneity and Wealth Normalization}
\label{subsec:homogeneity-normalization}
% --------------------------------------------------------------

\paragraph{Homothetic Scaling.}

With CRRA or Epstein-Zin preferences, the value function inherits the homogeneity of the utility function.

\begin{proposition}[Homogeneity of the Value Function]
\label{prop:homogeneity}
Under Epstein-Zin preferences with no labor income or binding constraints, the value function satisfies
\begin{equation}
V_t(\lambda W_t) = \lambda \cdot V_t(W_t) \quad \text{for all } \lambda > 0.
\label{eq:homogeneity}
\end{equation}
That is, $V_t$ is homogeneous of degree one in wealth.
\end{proposition}

\begin{proof}
By induction. At terminal date $T$ (if finite horizon), $V_T(W) = W$ or $V_T(W) = u(W)$ which is homogeneous. For the inductive step, suppose $V_{t+1}(\lambda W) = \lambda V_{t+1}(W)$. Then the Bellman equation at $t$:
\begin{align*}
V_t(\lambda W_t) &= \max_{c, \balpha} \left[ (1-\beta) c^{1-1/\theta} + \beta \left( \E_t[V_{t+1}((\lambda W_t - c) R^p)^{1-\gamma}] \right)^{\frac{1-1/\theta}{1-\gamma}} \right]^{\frac{1}{1-1/\theta}}.
\end{align*}
Substituting $c = \lambda \tilde{c}$ and using $V_{t+1}(\lambda(\tilde{W}_t - \tilde{c}) R^p) = \lambda V_{t+1}((\tilde{W}_t - \tilde{c}) R^p)$:
\begin{align*}
V_t(\lambda W_t) &= \lambda \max_{\tilde{c}, \balpha} \left[ (1-\beta) \tilde{c}^{1-1/\theta} + \beta \left( \E_t[V_{t+1}((\tilde{W}_t - \tilde{c}) R^p)^{1-\gamma}] \right)^{\frac{1-1/\theta}{1-\gamma}} \right]^{\frac{1}{1-1/\theta}} \\
&= \lambda \cdot V_t(W_t).
\end{align*}
The factor $\lambda$ pulls out by the homogeneity of the CES aggregator.
\end{proof}

\paragraph{The Normalized Value Function.}

Homogeneity implies we can write
\begin{equation}
V_t(W_t) = v_t \cdot W_t,
\label{eq:normalized-value}
\end{equation}
where $v_t$ is the \emph{value per unit wealth}---a ``normalized'' value function that depends on state variables other than wealth (investment opportunities, time to horizon, etc.) but not on wealth itself.

This normalization reduces the dimensionality of the problem: instead of solving for $V_t(W)$ over a wealth grid, we solve for the scalar $v_t$ (or a function of non-wealth states).

\paragraph{Normalized Consumption and Portfolio Rules.}

The optimal policies also normalize:
\begin{align}
c_t\opt &= m_t \cdot W_t, \label{eq:normalized-consumption} \\
\balpha_t\opt &= \balpha_t\opt(\text{states other than } W_t). \label{eq:normalized-portfolio}
\end{align}
The consumption-wealth ratio $m_t$ and portfolio weights $\balpha_t\opt$ are independent of the wealth level.

% --------------------------------------------------------------
\subsection{The Normalized Bellman Equation}
\label{subsec:normalized-bellman}
% --------------------------------------------------------------

Substituting the normalized form into the Bellman equation yields a recursion for $v_t$.

\paragraph{Derivation.}

Start with
\begin{equation}
v_t W_t = \max_{c_t, \balpha} \left[ (1-\beta) c_t^{1-1/\theta} + \beta \left( \E_t[(v_{t+1} W_{t+1})^{1-\gamma}] \right)^{\frac{1-1/\theta}{1-\gamma}} \right]^{\frac{1}{1-1/\theta}}.
\label{eq:bellman-substituted}
\end{equation}

Using $c_t = m_t W_t$ and $W_{t+1} = (1 - m_t) W_t R_{t+1}^p$:
\begin{align}
v_t W_t &= \left[ (1-\beta) (m_t W_t)^{1-1/\theta} + \beta \left( \E_t[(v_{t+1} (1-m_t) W_t R_{t+1}^p)^{1-\gamma}] \right)^{\frac{1-1/\theta}{1-\gamma}} \right]^{\frac{1}{1-1/\theta}} \notag \\
&= W_t \left[ (1-\beta) m_t^{1-1/\theta} + \beta (1-m_t)^{1-1/\theta} \left( \E_t[(v_{t+1} R_{t+1}^p)^{1-\gamma}] \right)^{\frac{1-1/\theta}{1-\gamma}} \right]^{\frac{1}{1-1/\theta}}.
\label{eq:bellman-factored}
\end{align}

Dividing by $W_t$:
\begin{equation}
\boxed{v_t = \max_{m_t, \balpha} \left[ (1-\beta) m_t^{1-1/\theta} + \beta (1-m_t)^{1-1/\theta} \left( \E_t[(v_{t+1} R_{t+1}^p)^{1-\gamma}] \right)^{\frac{1-1/\theta}{1-\gamma}} \right]^{\frac{1}{1-1/\theta}}.}
\label{eq:normalized-bellman}
\end{equation}

This is the \emph{normalized Bellman equation}. The state variable $W_t$ has been eliminated entirely.

\paragraph{The Certainty Equivalent Return (Revisited).}

Define
\begin{equation}
\tilde{R}_t^{CE}(\balpha) \equiv \left( \E_t\left[ (v_{t+1} R_{t+1}^p(\balpha))^{1-\gamma} \right] \right)^{\frac{1}{1-\gamma}}.
\label{eq:ce-return-normalized}
\end{equation}

This is the certainty equivalent of the ``value-weighted'' portfolio return. The normalized Bellman equation becomes
\begin{equation}
v_t = \max_{m_t, \balpha} \left[ (1-\beta) m_t^{1-1/\theta} + \beta (1-m_t)^{1-1/\theta} \left( \tilde{R}_t^{CE}(\balpha) \right)^{1-1/\theta} \right]^{\frac{1}{1-1/\theta}}.
\label{eq:bellman-with-ce}
\end{equation}

% --------------------------------------------------------------
\subsection{First-Order Conditions}
\label{subsec:foc}
% --------------------------------------------------------------

The normalized Bellman equation \eqref{eq:normalized-bellman} is optimized over the consumption rate $m_t$ and portfolio weights $\balpha$.

\paragraph{The Portfolio FOC.}

The portfolio weights enter only through $\tilde{R}_t^{CE}(\balpha)$. Maximizing over $\balpha$ is equivalent to maximizing the certainty equivalent:
\begin{equation}
\balpha\opt = \arg\max_{\balpha} \; \tilde{R}_t^{CE}(\balpha) = \arg\max_{\balpha} \; \E_t\left[ (v_{t+1} R_{t+1}^p(\balpha))^{1-\gamma} \right].
\label{eq:portfolio-foc-max}
\end{equation}

The first-order condition for an interior solution with respect to $\alpha_j$ is
\begin{equation}
\E_t\left[ (v_{t+1} R_{t+1}^p)^{-\gamma} \cdot v_{t+1} \cdot (R_{t+1}^j - R_{t+1}^k) \right] = 0 \quad \text{for all } j, k.
\label{eq:portfolio-foc}
\end{equation}

With a risk-free asset (return $R^f$), this simplifies to
\begin{equation}
\E_t\left[ (v_{t+1} R_{t+1}^p)^{-\gamma} \cdot v_{t+1} \cdot (R_{t+1}^j - R^f) \right] = 0 \quad \text{for all risky } j.
\label{eq:portfolio-foc-riskfree}
\end{equation}

\begin{calloutnote}[Interpretation of the Portfolio FOC]
Define the \emph{stochastic discount factor} (SDF):
\begin{equation}
M_{t+1} \equiv \beta \left( \frac{v_{t+1}(1-m_t) R_{t+1}^p}{v_t} \right)^{-1/\theta} \cdot \left( \frac{v_{t+1} R_{t+1}^p}{\tilde{R}_t^{CE}} \right)^{1/\theta - \gamma}.
\label{eq:sdf-ez}
\end{equation}
The portfolio FOC can be written as the asset pricing condition:
\begin{equation}
\E_t[M_{t+1} R_{t+1}^j] = 1 \quad \text{for all assets } j.
\label{eq:asset-pricing}
\end{equation}
This connects portfolio choice to the consumption-based asset pricing literature.
\end{calloutnote}

\paragraph{The Consumption FOC.}

Taking the first-order condition of \eqref{eq:bellman-with-ce} with respect to $m_t$:
\begin{align}
(1-\beta) m_t^{-1/\theta} &= \beta (1-m_t)^{-1/\theta} \left( \tilde{R}_t^{CE} \right)^{1-1/\theta}.
\label{eq:consumption-foc}
\end{align}

Rearranging:
\begin{equation}
\frac{m_t}{1-m_t} = \left( \frac{1-\beta}{\beta} \right)^{\theta} \left( \tilde{R}_t^{CE} \right)^{\theta - 1}.
\label{eq:consumption-foc-ratio}
\end{equation}

This confirms equation \eqref{eq:mpc-formula} from Section~\ref{subsec:nesting}.

\paragraph{Solving for the Normalized Value.}

Substituting the optimal $m_t\opt$ back into \eqref{eq:bellman-with-ce}, after algebra:
\begin{equation}
v_t = \left[ (1-\beta)^{\theta} + \beta^{\theta} \left( \tilde{R}_t^{CE} \right)^{\theta - 1} \right]^{\frac{1}{\theta - 1}} \cdot (1-\beta)^{\frac{\theta}{1-\theta}}.
\label{eq:v-solution}
\end{equation}

For the important case $\theta = 1$ (log utility over time), this simplifies dramatically:
\begin{equation}
v_t = (1-\beta) \exp\left( \frac{\beta}{1-\beta} \E_t[\log(v_{t+1} R_{t+1}^p)] \right).
\label{eq:v-log-utility}
\end{equation}

% --------------------------------------------------------------
\subsection{The Euler Equation}
\label{subsec:euler}
% --------------------------------------------------------------

The Euler equation characterizes intertemporal optimality by equating marginal utilities across time.

\paragraph{Derivation.}

The envelope condition for the Bellman equation gives
\begin{equation}
\frac{\partial V_t}{\partial W_t} = v_t.
\label{eq:envelope}
\end{equation}

The Euler equation for consumption is
\begin{equation}
(1-\beta) c_t^{-1/\theta} = \beta \E_t\left[ \frac{\partial V_{t+1}}{\partial W_{t+1}} \cdot R_{t+1}^p \right] \cdot \left( \frac{\mu_t}{\E_t[V_{t+1}^{1-\gamma}]^{1/(1-\gamma)}} \right)^{1/\theta - 1},
\label{eq:euler-raw}
\end{equation}
where the last term adjusts for the Epstein-Zin recursion.

\paragraph{The Consumption Euler Equation.}

After simplification, the Euler equation takes the form
\begin{equation}
\boxed{1 = \E_t\left[ \beta^{\theta} \left( \frac{c_{t+1}}{c_t} \right)^{-1/\theta} \left( \frac{v_{t+1} R_{t+1}^p}{\tilde{R}_t^{CE}} \right)^{1 - 1/\theta - (1-\gamma)} R_{t+1}^j \right]}
\label{eq:euler-ez}
\end{equation}
for any asset $j$ with return $R_{t+1}^j$.

\paragraph{Special Case: Time-Separable Utility.}

When $\gamma = 1/\theta$, the exponent $1 - 1/\theta - (1-\gamma) = 0$, and the Euler equation reduces to
\begin{equation}
1 = \E_t\left[ \beta \left( \frac{c_{t+1}}{c_t} \right)^{-\gamma} R_{t+1}^j \right],
\label{eq:euler-crra}
\end{equation}
the familiar CRRA Euler equation.

\paragraph{The Wealth Euler Equation.}

An equivalent formulation uses wealth growth instead of consumption growth. Since $c_t = m_t W_t$:
\begin{equation}
\frac{c_{t+1}}{c_t} = \frac{m_{t+1}}{m_t} \cdot \frac{W_{t+1}}{W_t} = \frac{m_{t+1}}{m_t} \cdot (1-m_t) R_{t+1}^p.
\label{eq:consumption-growth}
\end{equation}

In the i.i.d. case where $m_t = m$ is constant:
\begin{equation}
\frac{c_{t+1}}{c_t} = (1-m) R_{t+1}^p,
\label{eq:consumption-growth-iid}
\end{equation}
and the Euler equation depends only on the portfolio return.

% --------------------------------------------------------------
\subsection{Stationary Solutions}
\label{subsec:stationary}
% --------------------------------------------------------------

When the investment opportunity set is constant (i.i.d. returns), the problem admits a stationary solution.

\paragraph{The i.i.d. Case.}

Suppose returns $(R^1, \ldots, R^J)$ are i.i.d. over time with a fixed distribution. Then:
\begin{itemize}[nosep]
    \item The normalized value $v_t = v$ is constant
    \item The consumption rate $m_t = m$ is constant
    \item The portfolio weights $\balpha_t = \balpha$ are constant
\end{itemize}

\paragraph{Characterizing the Stationary Solution.}

The stationary normalized value satisfies
\begin{equation}
v = \left[ (1-\beta) m^{1-1/\theta} + \beta (1-m)^{1-1/\theta} \left( v \cdot \CE(R^p) \right)^{1-1/\theta} \right]^{\frac{1}{1-1/\theta}},
\label{eq:stationary-v}
\end{equation}
where $\CE(R^p) = (\E[(R^p)^{1-\gamma}])^{1/(1-\gamma)}$ is the certainty equivalent of the portfolio return.

The stationary consumption rate satisfies
\begin{equation}
\frac{m}{1-m} = \left( \frac{1-\beta}{\beta} \right)^{\theta} \left( v \cdot \CE(R^p) \right)^{\theta - 1}.
\label{eq:stationary-m}
\end{equation}

These two equations jointly determine $(v, m)$ given the optimal portfolio return distribution.

\begin{proposition}[Stationary Solution]
\label{prop:stationary}
With i.i.d. returns, the stationary normalized value is
\begin{equation}
v = \left( \frac{1-\beta}{1 - \beta \cdot \CE(R^{p,*})^{(\theta-1)/\theta}} \right)^{\theta},
\label{eq:v-stationary-closed}
\end{equation}
and the stationary consumption rate is
\begin{equation}
m = 1 - \beta^{\theta} \cdot \CE(R^{p,*})^{\theta - 1} \cdot (1-\beta)^{1-\theta}.
\label{eq:m-stationary-closed}
\end{equation}
These are well-defined when $\beta \cdot \CE(R^{p,*})^{(\theta-1)/\theta} < 1$.
\end{proposition}

\begin{calloutimportant}[The Complete Solution Algorithm]
For the i.i.d. case, the complete solution proceeds as:

\textbf{Step 1:} Solve for optimal portfolio weights $\balpha\opt$ by maximizing $\CE(R^p(\balpha)) = (\E[(R^p)^{1-\gamma}])^{1/(1-\gamma)}$. This is a static optimization problem.

\textbf{Step 2:} Compute $\CE(R^{p,*})$ at the optimal portfolio.

\textbf{Step 3:} Compute the stationary normalized value $v$ from \eqref{eq:v-stationary-closed}.

\textbf{Step 4:} Compute the consumption rate $m$ from \eqref{eq:m-stationary-closed}.

\textbf{Step 5:} The policy functions are $c_t = m W_t$ and $\balpha_t = \balpha\opt$.

When returns are not i.i.d. (e.g., predictable returns, stochastic volatility), Step 1 must be solved jointly with a recursion for $v_t$ as a function of the state variables describing investment opportunities.
\end{calloutimportant}



% Section 5: The KL Representation
% ==============================================================
% Section 5: The KL Representation
% ==============================================================

\section{The KL Representation}
\label{sec:kl-representation}

This section reveals a deep connection between portfolio choice and information theory. The certainty equivalent that governs risk preferences admits a variational representation involving the Kullback-Leibler divergence. This perspective unifies portfolio optimization with robust control, connects diversification to entropy, and provides bounds useful for asset pricing.

% --------------------------------------------------------------
\subsection{The Variational Representation}
\label{subsec:variational}
% --------------------------------------------------------------

Recall from Chapter~\ref{ch:information} that the classical generalized mean admits a Gibbs variational representation. We now apply this to the certainty equivalent.

\paragraph{The Gibbs Principle for Risk.}

The certainty equivalent of a random variable $X$ under probability measure $\mathbb{P}$ with risk aversion $\gamma$ is
\begin{equation}
\CE_{\gamma}(X) = \left( \E^{\mathbb{P}}[X^{1-\gamma}] \right)^{\frac{1}{1-\gamma}}.
\label{eq:ce-definition}
\end{equation}

When $\gamma > 1$ (the empirically relevant case), this can be written as a minimum over distorted probability measures.

\begin{proposition}[Variational Representation of Certainty Equivalent]
\label{prop:variational-ce}
For $\gamma > 1$ and $X > 0$:
\begin{equation}
\log \CE_{\gamma}(X) = \min_{\mathbb{Q} \ll \mathbb{P}} \left\{ \E^{\mathbb{Q}}[\log X] + \frac{1}{\gamma - 1} \KL{\mathbb{Q}}{\mathbb{P}} \right\}.
\label{eq:variational-ce}
\end{equation}
The minimum is attained at the \emph{distorted measure}
\begin{equation}
\frac{d\mathbb{Q}^*}{d\mathbb{P}} = \frac{X^{1-\gamma}}{\E^{\mathbb{P}}[X^{1-\gamma}]}.
\label{eq:distorted-measure}
\end{equation}
\end{proposition}

\begin{proof}
This follows from the Gibbs variational principle applied to the classical mean with power $\rho = 1 - \gamma < 0$. The sign of the KL term reflects that for $\rho < 0$, the mean is below the geometric mean, achieved by tilting toward low outcomes.
\end{proof}

\paragraph{Interpretation.}

The variational representation \eqref{eq:variational-ce} says: the certainty equivalent equals the expected log payoff under the \emph{worst-case} probability measure $\mathbb{Q}$, where ``worst-case'' is penalized by KL divergence from the true measure $\mathbb{P}$.

\begin{itemize}[nosep]
    \item The investor acts as if nature chooses $\mathbb{Q}$ adversarially
    \item Nature wants to minimize expected payoff but faces an entropy cost
    \item Higher risk aversion $\gamma$ $\Rightarrow$ lower penalty $(1/(\gamma-1))$ $\Rightarrow$ more distortion toward bad states
\end{itemize}

% --------------------------------------------------------------
\subsection{Connection to Robust Control}
\label{subsec:robust-control}
% --------------------------------------------------------------

The variational representation connects naturally to the robust control framework of Hansen and Sargent (2001, 2008).

\paragraph{Model Uncertainty.}

Hansen and Sargent consider an agent who distrusts the baseline probability model $\mathbb{P}$. The agent evaluates decisions under a worst-case model $\mathbb{Q}$ from a set of plausible alternatives:
\begin{equation}
V = \min_{\mathbb{Q}: \KL{\mathbb{Q}}{\mathbb{P}} \leq \eta} \E^{\mathbb{Q}}[U].
\label{eq:robust-minmax}
\end{equation}

The Lagrangian relaxation of \eqref{eq:robust-minmax} is
\begin{equation}
V = \min_{\mathbb{Q}} \left\{ \E^{\mathbb{Q}}[U] + \theta \KL{\mathbb{Q}}{\mathbb{P}} \right\},
\label{eq:robust-lagrangian}
\end{equation}
where $\theta > 0$ is the Lagrange multiplier on the entropy constraint.

\paragraph{Observational Equivalence.}

Comparing \eqref{eq:robust-lagrangian} with \eqref{eq:variational-ce} (taking $U = \log X$):
\begin{equation}
\theta = \frac{1}{\gamma - 1}.
\label{eq:theta-gamma}
\end{equation}

\begin{calloutimportant}[Risk Aversion vs. Robustness]
The variational representation reveals an \emph{observational equivalence}:
\begin{center}
\begin{tabular}{@{}lcc@{}}
\toprule
\textbf{Interpretation} & \textbf{Parameter} & \textbf{Mechanism} \\
\midrule
Risk aversion & $\gamma$ & Curvature of utility \\
Robustness/ambiguity & $\theta = 1/(\gamma-1)$ & Distrust of probability model \\
\bottomrule
\end{tabular}
\end{center}

An agent with high risk aversion ($\gamma = 10$) behaves identically to an agent with log utility who distrusts the model with penalty $\theta = 1/9$.

This equivalence has profound implications: observed ``risk averse'' behavior may reflect either genuine utility curvature or model uncertainty. Disentangling these requires auxiliary data (e.g., behavior under known probabilities).
\end{calloutimportant}

\paragraph{The Hansen-Sargent Multiplier.}

Hansen and Sargent (2008) develop a full dynamic theory where the agent solves
\begin{equation}
V_t = \max_{c, \balpha} \min_{\mathbb{Q}} \left\{ u(c) + \beta \E_t^{\mathbb{Q}}[V_{t+1}] + \beta \theta \KL{\mathbb{Q}_t}{\mathbb{P}_t} \right\}.
\label{eq:hs-bellman}
\end{equation}

The inner minimization distorts beliefs about next-period states; the outer maximization chooses consumption and portfolio. The solution has the same structure as Epstein-Zin with risk aversion $\gamma = 1 + 1/\theta$, but the \emph{interpretation} differs: the agent has unit risk aversion but fears model misspecification.

\begin{calloutnote}[Detection Error Probabilities]
Hansen and Sargent calibrate $\theta$ using detection error probabilities: how much data would be needed to statistically distinguish $\mathbb{Q}^*$ from $\mathbb{P}$? If the worst-case model is ``close'' to the baseline (high KL cost for small distortions), robustness concerns are mild. Their calibrations suggest $\theta$ values corresponding to $\gamma \approx 5$--$20$.
\end{calloutnote}

% --------------------------------------------------------------
\subsection{Entropy and Diversification}
\label{subsec:entropy-diversification}
% --------------------------------------------------------------

The KL representation also illuminates the diversification motive in portfolio choice.

\paragraph{Portfolio Entropy.}

Consider portfolio weights $\balpha = (\alpha_1, \ldots, \alpha_J)$ with $\alpha_j \geq 0$ and $\sum_j \alpha_j = 1$. The \emph{entropy} of the portfolio is
\begin{equation}
H(\balpha) = -\sum_{j=1}^{J} \alpha_j \log \alpha_j.
\label{eq:portfolio-entropy}
\end{equation}

Entropy is maximized by the equal-weighted portfolio $\alpha_j = 1/J$, which achieves $H = \log J$. Concentrated portfolios have low entropy.

\paragraph{Diversification as Entropy Maximization.}

Suppose asset returns are exchangeable (symmetric distribution). Then the optimal portfolio maximizes a trade-off:
\begin{equation}
\max_{\balpha} \left\{ \E[\log R^p(\balpha)] - \frac{1}{\gamma - 1} H(\balpha) \right\} \quad \text{(for } \gamma < 1 \text{)}
\label{eq:entropy-tradeoff}
\end{equation}
or the analogous minimization for $\gamma > 1$.

The entropy term penalizes concentration. Even with identical expected returns, a risk-averse investor diversifies because concentration increases exposure to idiosyncratic outcomes.

\paragraph{The Maximum Entropy Portfolio.}

In the absence of informative signals about expected returns, the principle of maximum entropy suggests the equal-weighted portfolio:
\begin{equation}
\balpha^{ME} = \left( \frac{1}{J}, \ldots, \frac{1}{J} \right).
\label{eq:max-entropy-portfolio}
\end{equation}

DeMiguel, Garlappi, and Uppal (2009) document that this naive ``$1/N$'' portfolio often outperforms optimized portfolios out-of-sample, precisely because estimation error in expected returns can overwhelm the gains from optimization.

% --------------------------------------------------------------
\subsection{Bounds and Asset Pricing}
\label{subsec:bounds-asset-pricing}
% --------------------------------------------------------------

The variational representation provides useful bounds for asset pricing and connects to the Hansen-Jagannathan framework.

\paragraph{The Entropy Bound.}

From \eqref{eq:variational-ce}, for any probability measure $\mathbb{Q}$:
\begin{equation}
\log \CE_{\gamma}(X) \leq \E^{\mathbb{Q}}[\log X] + \frac{1}{\gamma - 1} \KL{\mathbb{Q}}{\mathbb{P}}.
\label{eq:entropy-bound}
\end{equation}

Choosing $\mathbb{Q} = \mathbb{P}$ (no distortion):
\begin{equation}
\log \CE_{\gamma}(X) \leq \E^{\mathbb{P}}[\log X].
\label{eq:ce-log-bound}
\end{equation}

This confirms that for $\gamma > 1$, the certainty equivalent lies below the geometric mean.

\paragraph{The Stochastic Discount Factor and Entropy.}

Let $M$ be the stochastic discount factor pricing all assets: $\E[MR^j] = 1$. Alvarez and Jermann (2005) show that the entropy of the SDF,
\begin{equation}
L(M) = \log \E[M] - \E[\log M],
\label{eq:sdf-entropy}
\end{equation}
provides a lower bound on the maximum expected log return:
\begin{equation}
\max_j \E[\log R^j] \leq L(M).
\label{eq:alvarez-jermann}
\end{equation}

This ``entropy bound'' complements the Hansen-Jagannathan volatility bound. While the HJ bound constrains $\sigma(M)/\E[M]$, the entropy bound constrains the ``size'' of multiplicative SDF fluctuations.

\paragraph{Implications for the Equity Premium.}

The equity premium puzzle can be restated in entropy terms. If the SDF is $M = \beta(c_{t+1}/c_t)^{-\gamma}$, then
\begin{equation}
L(M) \approx \frac{\gamma^2}{2} \sigma^2(\Delta \log c).
\label{eq:sdf-entropy-crra}
\end{equation}

With consumption growth volatility $\sigma(\Delta \log c) \approx 2\%$ and equity premium $\E[\log R^e - \log R^f] \approx 6\%$, the bound requires $\gamma \geq 30$---the puzzle in entropy form.

\begin{calloutnote}[Bansal-Yaron Long-Run Risk]
Bansal and Yaron (2004) resolve the entropy bound by introducing persistent consumption growth shocks. With long-run risk, small shocks to expected growth cumulate over time, generating large entropy in the SDF without requiring extreme risk aversion. The Epstein-Zin preference for early resolution of uncertainty (Section~\ref{subsec:resolution-timing}) amplifies the pricing of long-run risk.
\end{calloutnote}

% --------------------------------------------------------------
\subsection{The Cost of Suboptimal Portfolios}
\label{subsec:cost-suboptimal}
% --------------------------------------------------------------

The KL representation quantifies the welfare cost of deviating from the optimal portfolio.

\paragraph{Welfare Loss.}

Let $\balpha\opt$ be the optimal portfolio and $\balpha$ any suboptimal portfolio. The welfare loss, measured in certainty-equivalent units, is
\begin{equation}
\Delta \CE = \CE_{\gamma}(R^{p,*}) - \CE_{\gamma}(R^p(\balpha)).
\label{eq:welfare-loss}
\end{equation}

Using the variational representation, this can be bounded in terms of the KL divergence between the distorted measures induced by the two portfolios.

\paragraph{Proportional Loss.}

For small deviations from optimality, the proportional loss is approximately
\begin{equation}
\frac{\Delta \CE}{\CE\opt} \approx \frac{\gamma}{2} \cdot \frac{\Var(R^p - R^{p,*})}{(\CE\opt)^2}.
\label{eq:proportional-loss}
\end{equation}

The loss is quadratic in the deviation and proportional to risk aversion. This justifies the common approximation that small portfolio mistakes have second-order welfare effects---though with high $\gamma$, even modest deviations can be costly.

\paragraph{The Value of Information.}

The KL framework also quantifies the value of information about returns. If an investor can acquire a signal $s$ about returns at cost $\kappa$, the value of information is
\begin{equation}
\text{VoI} = \E_s\left[ \CE_{\gamma}(R^{p,*}(s)) \right] - \CE_{\gamma}(R^{p,*}(\text{no signal})) - \kappa.
\label{eq:value-of-info}
\end{equation}

Under rational inattention (Sims 2003, Ma\'{c}kowiak and Wiederholt 2009), the optimal attention allocation balances the marginal value of information against its Shannon entropy cost---another KL-based trade-off.

\begin{calloutimportant}[Summary: KL Divergence in Portfolio Choice]
The KL divergence appears in multiple guises:
\begin{enumerate}[nosep]
    \item \textbf{Variational representation}: $\CE_{\gamma}$ as worst-case expectation with entropy penalty
    \item \textbf{Robustness}: Equivalent to model uncertainty aversion (Hansen-Sargent)
    \item \textbf{Diversification}: Entropy of portfolio weights as a diversification measure
    \item \textbf{Asset pricing bounds}: SDF entropy bounds expected returns
    \item \textbf{Welfare costs}: KL between optimal and suboptimal portfolio distributions
    \item \textbf{Information value}: Shannon cost of acquiring return signals
\end{enumerate}
This unifying structure connects portfolio theory to information theory, statistical mechanics, and robust control.
\end{calloutimportant}


% \input{sec3_consumption_portfolio}
% \input{sec4_value_function}
% \input{sec5_kl_representation}